\section*{Executive summary}

Severe large-vessel occlusion stroke is best treated by a combination of thrombolysis (IVT) and thrombectomy (MT). Most patients first attend an IVT-only centre, and require transfer for thrombectomy. Those patients requiring transfer have a lower chance of receiving thrombectomy, and those that receive thrombectomy receive it about 110 minutes later than patients who directly attend an MT-capable unit. The delay is due to a combination of extra travel time, transferring between units, and system-related delays in initiating transfer between stroke units. Improving access to MT, especially direct access, is expected to improve outcomes in the most severe strokes.

Introducing thrombectomy at Taunton and Worcester would give about three more large vessel occlusion stroke patients in every 100 treated there a good recovery (independent living at 3–6 months), with Taunton showing the greater benefit. That gain is uneven across patients. Those who go straight to the new centres gain five to eight extra good outcomes per 100 treated. Those transferred in from thrombolysis-only hospitals gain around one. To get the most benefit, ambulance services should take suitable patients straight to the new centres where possible.